\documentclass[a4paper,12pt]{article}

%need to format subitems

\setcounter{tocdepth}{2} % show only 2 levels in toc

\begin{document}

\tableofcontents
\newpage

\section{Capitalism - A Global History, Sven Beckert}
\subsection{Chapter 3 - The Great Connectiong, 1450-1650}
\begin{itemize}
\item Prior to this, traded goods did not penetrate local markets, but they allowed for capital accumulation, which was not yet possible in agriculture or industry.

\item We begin to see merchants moving to distant places, rather than just trading with them. (does this not contradict the innovation of chapter 1? large trading networks controlled by merchants in place?)

\item In the 15th century, merchants began to invest in agriculture and manufacturing.

\item prior to this, there needed to be a 'great connecting'

\subitem powered by a cooperation between merchants and state

\item this great connecting was centered on europe because of their control of the atlantic - great divergence?

\subitem why europe? because europe was relatively weak and needed to over come this

\subitem first: ``European capital owners forged new islands of capital because they had great incentive to do so.'' - they had to circumvent muslem middlemen, thus the search for alternative routes to asia - columbus, cape of good hope etc.

\subitem second: European states were facing crises of Feudalism. Especially labor shortage.

\subitem third: European states needed funding for their wars.

\subitem ``The relative weakness of European states and capitalists on the global stage drove them closer to one another and toward a consensus on projecting their power and capital into novel spaces and constructing entirely new islands of capital.''

\item European capital set up productive outposts in the New World. Slavery powered the colonies, and itself became quite profitable. Labor itself first became a commodity here in the form of enslaved people.

\item European cities such as Amsterdam became ruled by the logic of capital: the market first penetrated everyday life here. But it would be small scale for a while.

\item The great connecting created institutions:

\subitem Joint stock companies

\subitem Futures trading and stock exchanges

\subitem The family - the division of labor within the household

\subitem The state - backed the other institutions; upheld property rights; involved the military in trade and capitalism.

\item With the increased involvement of states in trade, merchants also began to be more involved in politics: ``rulers everywhere became more dependent on the fortunes and whims of the capitalists whose expansion they had enabled.''

\item Back to the great divergence: merchants in china faced a more sophisticated and inflexible state that limited their influence

\item The example of the Fuggers: combined state and merchant power. Built the first factories (i.e. storehouses), received mineral rights from states, etc.

\item Sometimes the merchants took on statelike functions themselves, as with the East India companies. 

\item War capitalism: during this period, accumulation was largely driven by ``coercive expansion of European nodes of capital''

\subitem unlike tributary economies, here the redistributed wealth went to merchants rather than rulers.

\item As Europeans began to dominate trade, they undermined local merchants and existing muslim trade routes. They used force to create monopolies for themselves.

\end{itemize}

\subsection{Chapter 4 - Transforming the Countryside, 1550-1750}

\begin{itemize}

\item Barbados became ``one of the world's first recognizably capitalist societies''

\item land becomes alienable in 17th century

\item slaves become a form of fungible collateral for loans (market in labor)

\item different form of labor control
	\subitem indigenous slavery
	\subitem indentured workers
	\subitem serfdom
	\subitem african slavery

\end{itemize}

\end{document}
